\section{Theory: non-identifiability and the necessity of abstention}
\label{sec:theory}

\todo{Tighten proof prose; confirm notation matches Sec.~4--5; add related-work pointers for
partial identification. Content below is locked to csc/THEORY.md.}

\paragraph{Setup.}
Let $Z$ be a fixed representation of an EEG epoch and $Y$ the label. The source law is $P_S(Z,Y)$; a target
provides samples of $Q_Z$ (the unlabeled target marginal), and in the minimal-label regimes of
Sec.~\ref{sec:routeB3method} a small number of target labels. A \emph{concept shift} is a change in the
posterior, $Q(Y\mid Z)\neq P_S(Y\mid Z)$.

\subsection{A $Z$-only certificate cannot identify concept shift}

\begin{proposition}[No $Z$-only concept identification]
\label{prop:impossible}
Fix any target marginal $Q_Z$. There exist two joint target laws $Q_0,Q_1$ with
$Q_0(Z)=Q_1(Z)=Q_Z$ but $Q_0(Y\mid Z)=P_S(Y\mid Z)$ and $Q_1(Y\mid Z)\neq P_S(Y\mid Z)$. Any certificate
that is a function of the unlabeled target (a function of $Q_Z$) returns the same value on $Q_0$ and $Q_1$.
Hence no $Z$-only certificate can decide ``$P(Y\mid Z)$ unchanged'' vs.\ ``concept shift'' --- for
\emph{any} observed marginal, not only when $Q_Z=P_S(Z)$.
\end{proposition}

\begin{proof}
Pick any $Q_Z$ with full support on the source $Z$-support and set $Q_0(Y\mid Z)=P_S(Y\mid Z)$. For $Q_1$,
perturb the posterior on a positive-measure set (e.g.\ swap two classes' posteriors on a region) while
keeping the same $Q_Z$; both are valid joints with marginal $Q_Z$. A $Z$-only $f$ sees $Q_Z$ in both cases,
so $f(Q_0)=f(Q_1)$ and is wrong on at least one. \qed
\end{proof}

\paragraph{Consequence.}
``The target marginal landed in a concept direction'' is, on its own, \emph{not} identification of concept
shift; it is only \emph{evidence} under an extra, explicit assumption that ties what the source showed to
what the target is doing.

\begin{assumption}[atlas transportability + signature separability]
\label{ass:T}
(i) The target shift is generated by the same mechanism family the source spans --- its marginal signature
lies in the span of the source covariate/concept/label signatures (no novel mechanism); and (ii) those
signature subspaces are separable (a covariate-only move and a boundary-moving concept change leave
distinguishable marginal traces).
\end{assumption}

Under Assumption~\ref{ass:T} a marginal certificate is a \emph{conditional, abstention-equipped partial
identification} statement, never distribution-free identification --- and it must abstain whenever the
target leaves the atlas (Prop.~\ref{prop:impossible}'s ``novel direction'' branch).

\subsection{The shift taxonomy and the three-state output}

\begin{table}[t]
\centering
\caption{Shift taxonomy. Only a $Z$-visible \emph{and} source-anchored trace admits a (conditional)
certificate; pure conditional and label shift are $Z$-only-unidentifiable and must abstain.}
\label{tab:taxonomy}
\begin{tabular}{@{}llll@{}}
\toprule
shift & $P(Z)$ & $P(Y\mid Z)$ & marginal-only certificate \\
\midrule
covariate                 & changes (nuisance) & invariant & \stcov{} \\
boundary-coupled concept  & changes            & changes   & \stsus{} (under Assumption~\ref{ass:T}) \\
pure conditional concept  & invariant          & changes   & \stunid{} (Prop.~\ref{prop:impossible}) \\
label (target) shift      & changes            & changes   & \stunid{} \\
\bottomrule
\end{tabular}
\end{table}

Table~\ref{tab:taxonomy} forces a \textbf{three-state} $Z$-only output --- \stcov{} / \stsus{} / \stunid{}
--- never a binary detector. \stcov{} is a \emph{compatibility} claim (source boundary stable in the moved
subspace), not an adaptation guarantee. The label-shift row is the sharpest trap: a change in $P(Y)$ with
$P(Z\mid Y)$ fixed moves the pooled mean along the class-mean subspace and mimics concept movement, so the
certifier carries an explicit label subspace and label gate and abstains there.

\begin{figure}[t]
\centering
\includegraphics[width=0.86\linewidth]{fig1_taxonomy_boundary.png}
\caption{Shift taxonomy and the identifiability/abstention boundary: $Z$-visible, source-anchored traces
admit a conditional certificate; invisible/label-confounded/out-of-atlas shifts fall in the abstention
region (Prop.~\ref{prop:impossible}).}
\label{fig:taxonomy}
\end{figure}

\paragraph{What is proven vs.\ assumed (honesty).}
\emph{Proven} (Prop.~\ref{prop:impossible}): a $Z$-only certificate cannot identify concept shift for any
marginal. \emph{Assumed/open}: Assumption~\ref{ass:T} (transportability + separability) is untestable from
the target marginal alone; a marginal \stsus{} is therefore source-anchored suspicion under an explicit
assumption. This is precisely the gap that motivates supplying a small amount of \emph{additional information} in
Sec.~\ref{sec:routeB3method}.
